\section{GPT-3：Meta-Learning 与 Few-Shot 学习曲线}

GPT-3 进一步把 examples 放进 context，使模型在不更新参数时根据局部 demonstration 调整输出。课堂称之为 language-model meta-learning：outer loop 是预训练，inner-like behavior 发生在 context 中。这个说法是功能类比，不必假设模型真的执行梯度下降。

\subsection{Meta-learning interface}

本节先定义推理时的任务接口。\term{few-shot} 指 prompt 内提供少量 input-output examples；\term{in-context learning} 指模型仅根据这些上下文改变行为。对参数 $\theta$，预测仍是 $p_\theta(y\mid x,\text{demonstrations})$，$\theta$ 不变。

\lecturefigure{slide-18-gpt3-metalearning.jpg}{GPT-3 将 task description 和 demonstrations 放入同一 context。}{官方视频 18:15；Brown et al. 2020。}

\begin{knowledgebox}{读图：训练与推理的两层结构}
Pretraining 在大量任务样式上更新参数；inference prompt 提供临时任务定义。模型可能通过 pattern completion、retrieval-like matching 或 learned algorithm 实现 behavior。无论内部机制如何，工程上都需测试 demonstration order、label words、format 和 context length。
\end{knowledgebox}

\subsection{Few-shot arithmetic}

本节用 arithmetic curve 检查 scale 与 demonstrations 的交互。图中比较 zero/one/few-shot 与 model size；某些位数任务在规模增大后显著改善，但不同难度曲线不同，说明能力不是一个统一 scalar。

\lecturefigure{slide-19-few-shot-arithmetic.jpg}{GPT-3 arithmetic 曲线展示 scale、shots 与任务难度的交互。}{官方视频 19:23。}

\begin{knowledgebox}{读图：先看每条曲线代表什么}
横轴通常是 model size，纵轴是 accuracy；颜色或线型区分位数/shot setting。重点不是“更大总是更好”，而是某些任务斜率更陡、某些接近随机。算术正确还可能受 tokenization 与训练数据格式影响，不能直接等同通用数学推理。
\end{knowledgebox}

\subsection{Word unscrambling}

Unscrambling 要求模型从字符或词片段恢复单词。它比常规 language completion 更接近 algorithmic mapping，可观察 demonstrations 是否让模型推断规则。

\lecturefigure{slide-20-word-unscrambling.jpg}{Word unscrambling 用 few-shot examples 测试规则归纳。}{官方视频 20:14。}

\begin{knowledgebox}{读图：关注 task format sensitivity}
曲线提高说明更大模型更能利用 context examples，但字符分词、大小写、空格与答案格式都会改变难度。若只换 label token 就显著下降，说明模型学到的可能是脆弱 pattern，而不是稳定 algorithm。
\end{knowledgebox}

\subsection{General few-shot learning}

进一步把多任务表现与 scale 放在同一汇总图中。Average 可以显示总体趋势，也会隐藏任务间差异、负迁移和 benchmark saturation，因此需要回到 per-task evidence。

\lecturefigure{slide-21-general-few-shot.jpg}{GPT-3 多任务曲线把 zero/one/few-shot 与规模化趋势放在共同坐标中。}{官方视频 20:36。}

\begin{knowledgebox}{读图：平均线下面是任务分布}
比较 zero-shot、one-shot、few-shot 的间距，再看其是否随规模扩大。Few-shot 提升若只来自少数任务，平均分会误导；同时 benchmark 可能被 pretraining contamination 抬高。可靠结论需要 per-task breakdown、confidence 与 held-out data governance。
\end{knowledgebox}

\teachervoice{讲者把 GPT-3 的关键变化描述为“language model metalearning”。课堂提示是 interface：任务被放进 context，而不是为每项任务重新训练模型。这个接口后来很重要，但当时曲线已经显示它对任务和格式并不均匀。}

\subsection{本章小结}

GPT-3 让 demonstrations 成为推理时输入，形成 in-context interface。Scale 提高了使用上下文的能力，却没有消除格式敏感、任务差异与 contamination 风险。

